% ── Part I, Chapter 4 ────────────────────────────────────────
% Correspondence — propagation, each recipient reconstructing
% differently, the text fragmenting into local interpretations
% Guyon, France, 1680s

\chapter{Correspondence}

\strandmarker{Guyon}{France, 1680s}

\lettrine{A}{} letter is not the thing it describes.

She knows this and writes anyway, because the alternative
is silence, and silence does not propagate. What she has
found cannot remain only interior. Not because she requires
witnesses---she has examined this motive carefully and found
it mostly absent---but because the discovery, if it is a
discovery, implies something about the nature of the process
that does not restrict it to a single practitioner. If the
system can reach this configuration, it can reach it more
than once. Other systems, differently situated, differently
constrained, might find their way to the same place by
different paths.

She writes to describe the path, knowing that the description
is not the path.

The letters leave her hands and enter a different system
entirely: the system of reading, of interpretation, of
reconstruction from marks on a surface back into something
that can move inside a mind. At each stage of this
reconstruction, choices are made that she did not make and
cannot control. The reader brings a framework. The framework
selects from what the text offers. What is selected is not
the text but the text as it appears within that framework,
which is a different thing.

She cannot follow the letters to correct the selection.

\section{}

Fénelon reads her differently than she writes.

This is not a criticism. He reads with more precision than
almost anyone she has encountered, with a genuine effort
to follow the structure of what she is saying rather than
to immediately translate it into something already familiar.
He pauses at the places she intends to be paused at. He
returns to earlier passages when later ones require it.
He reads, in other words, as though the text is trying to
do something specific and his task is to understand what
that is rather than to assess whether it succeeds by his
own criteria.

And yet.

What arrives at him is not quite what she sent.

She can see this in his responses, which are careful and
generous and which reformulate her positions with a
delicacy that preserves the shape while smoothing what
was intentionally unsmooth. The places in her writing
that do not resolve---that are left open because closure
would be false---arrive at him as places that require
further work, that could, with sufficient refinement, be
brought to a stable formulation. He offers the refinement.
The refinement is elegant and wrong in a way that is
difficult to name, because it is wrong not in its content
but in its confidence.

She writes back with the unsmooth parts restored.

He reads this too, with the same care, and smooths them
again, because smoothing is what his system does with
what it cannot fully resolve, and he cannot fully resolve
the deliberate roughness any more than he can stop
reading carefully.

This is not failure. It is a property of translation.

\section{}

Others read her with less care and more certainty.

A letter arrives at a woman in Lyon who has been looking
for permission to trust what she has already been
experiencing, and finds it, and the experience expands
to fill the permission, and what she writes back to her
own correspondents is not what the letter said but what
the letter allowed. The expansion is genuine. The
experience is real. It has, however, migrated from the
letter's actual content to the letter's emotional
register, and these are not the same thing.

A letter arrives at a priest in Grenoble who reads it as
a system of practices, a method, a sequence of steps that
can be taught and followed. He extracts the steps. They
are not quite there in the original, but they are close
enough to a certain reading of the original, and he is
experienced at extracting steps, and so steps are
extracted. He teaches them. His students follow them.
Some of his students report the same experiences she
describes. Others do not, and conclude that they have
followed the steps incorrectly.

The steps were not the point.

She watches these receptions accumulate and feels something
she does not name as concern, because concern implies
that she expected otherwise. She did not expect otherwise.
She knew, before the first letter left her hand, that
the description is not the path. She wrote anyway because
the alternative was silence.

\section{}

The letters also travel to people she did not address.

They are copied, partially, inaccurately, without context.
Phrases are extracted and recombined with phrases from
other sources. A sentence that was doing careful work
in its original position is removed from that position
and placed somewhere else, where it does different work
or no work or actively misleading work. She cannot track
this. No one could track this. The text has entered a
system that does not preserve the conditions of its
production.

What circulates is increasingly a reconstruction of a
reconstruction, each iteration moving further from what
she wrote and closer to what the circulating system
needs her to have written.

This is how ideas travel.

Not as themselves but as the versions of themselves that
survive the passage through successive frameworks, each
one selecting for what it can use and discarding the
rest. What arrives at any given point in the network
is a composite: the original, compressed and distorted
by everything it has passed through, carrying traces
of each reconstruction in its current form.

She cannot prevent this.

She can only write more letters, more carefully, with
the difficult parts left difficult, and accept that
each letter will begin the same process again: passing
through hands that will do what hands do with what they
cannot fully hold, which is to reshape it into something
holdable.

The difficult parts matter most.

They are also the first to be smoothed away.

\section{}

Fénelon visits in the spring.

They speak for several hours, which is different from
the letters in ways she had expected and in ways she
had not. The letters allow precision at the cost of
presence. The conversation allows presence at the cost
of precision: she can see when he follows and when he
does not, can adjust in real time, can leave a thought
incomplete and watch whether he completes it in the
direction she intends or in a direction she did not.

He completes it, mostly, in directions she did not.

But he also---and this she had not anticipated---says
things that adjust her own understanding. Not because
he is further along than she is, but because the act
of being heard inaccurately clarifies what she meant
to say. His misreadings are diagnostic. They show her
where the writing is unclear not because the ideas
are unclear but because the writing has not yet found
the form the ideas require.

She makes notes after he leaves.

Not of what he said, but of where he went wrong, which
is a map of where the work still needs to happen.

The letters continue.

She writes into the gap between what she can say and
what she means to say, knowing the gap will not close,
finding in the attempt to close it something that is
not closure but is not nothing.

